% ---- III. Problem Formulation. Source: JAMA "Markov Decision Process Framework", expanded with
% capability tiers, subtype marginalization, in-hospital intervals and unrestricted transfers;
% eMethods 1 equations moved into the main text.
\section{Problem Formulation}
\label{sec:problem}

We model prehospital stroke triage as a Markov decision process (MDP)~\cite{bellman1957,kochenderfer2022} defined by a state space, an action space, a transition model, and a reward function. The objective is a policy, a mapping from states to actions, that maximizes expected reward. Here the horizon is short, an initial transport and at most one inter-facility transfer, and the reward is a single terminal quantity: the probability of a good outcome once the patient's care pathway is complete.

\subsection{Hospital capability tiers}
Hospitals are classified by what they can do for a stroke patient on arrival, not by certification label alone. An \emph{EVT-capable center} performs endovascular thrombectomy on site; in the primary analysis this means a Joint Commission Comprehensive or Thrombectomy-Capable Stroke Center, a DNV Comprehensive Stroke Center, or a hospital designated for thrombectomy by its county EMS agency. A \emph{thrombolysis-capable center} administers intravenous thrombolysis but transfers patients for EVT (Primary, Advanced Primary and Acute Stroke Ready certification). A \emph{non-stroke-center hospital} has a 24-hour emergency department and no stroke certification; we assume it neither administers thrombolysis nor performs EVT and transfers for both. Section~\ref{sec:setup} describes how designations were verified, and Section~\ref{sec:results} reports a sensitivity analysis in which hospitals known to perform thrombectomy without certification are treated as EVT-capable.

Fig.~\ref{fig:mdp} summarizes the states, actions, information flow and timing of the model.

\begin{figure*}[t]
\centering
\includegraphics[width=0.9\textwidth]{fig_influence_standalone}
\caption{Influence diagram of the routing decision. The initial state $S_0$ holds the pickup location, the time since onset and the stroke subtype $\theta$, which is drawn from the population prior and hidden from the crew; the first decision $A_0$ observes location and time but not $\theta$. Arrival at $h_1$ advances the clock by the road travel time and imaging reveals $\theta$. The second decision $A_1$, made with $\theta$ known, is to treat at $h_1$ or transfer once to $h_2$. The reward is paid once, at the terminal state, as the probability of an excellent outcome given the subtype and the needle and puncture times the completed pathway implies (Table~\ref{tab:intervals}); the first leg earns no reward. For hemorrhage and stroke mimics the outcome does not depend on time.}
\label{fig:mdp}
\end{figure*}

\subsection{State space}
A state $s = (\ell, t, k, \sigma)$ comprises the patient's location $\ell$ (a pickup coordinate or a hospital), minutes since symptom onset $t$, whether the subtype is known $k \in \{\text{unknown}, \text{known}\}$, and the subtype $\sigma \in \{\text{LVO}, \text{nLVO}, \text{ICH}, \text{mimic}\}$. Pickup locations are sampled from high-resolution population density data~\cite{meta2019}, and onset-to-pickup time from 30 to 270 minutes, reflecting typical ranges to EMS contact~\cite{eddelien2021}. The subtype is drawn from the published prevalence of each category among patients with a positive prehospital LVO screen (Los Angeles Motor Scale $\geq 4$)~\cite{holodinsky2018}: $p_{\text{LVO}} = 0.454$, $p_{\text{nLVO}} = 0.109$, $p_{\text{ICH}} = 0.345$, $p_{\text{mimic}} = 0.092$. The planner never observes the drawn subtype while $k = \text{unknown}$; it reasons over the prior (Section~\ref{sec:method}). The subtype becomes known on arrival at any hospital, when imaging is assumed to occur.

\subsection{Action space}
Actions are discrete and location-dependent, and never route a patient to a lower capability tier. From the field, the patient may be transported to any hospital the road network reaches. From a non-stroke-center hospital the options are transfer to a thrombolysis-capable or EVT-capable center, or stay. From a thrombolysis-capable center, transfer to an EVT-capable center or stay. At an EVT-capable center the only action is to stay.

\subsection{Transition model}
A routing action moves the patient to the chosen hospital and advances the clock by the road travel time $\tau(\ell, \ell')$ obtained from the OpenRouteService routing engine on OpenStreetMap data; a transfer out of a hospital additionally incurs the door-in-door-out interval $T_{\text{DIDO}}$ before departure. On arrival the subtype becomes known. Transitions are deterministic given the action; all uncertainty enters through the subtype.

\subsection{In-hospital intervals}
Treatment is not instantaneous on arrival. We model three intervals, one value per interval for all hospitals of a tier, taken from national benchmarks (Table~\ref{tab:intervals}): door-to-needle $T_{\text{DTN}}$ for thrombolysis; door-to-puncture $T_{\text{DTP}}$ for EVT after direct arrival, and a shorter $T_{\text{DTP}}^{\text{tr}}$ for a transferred-in patient who arrives already imaged; and door-in-door-out $T_{\text{DIDO}}$ at the sending hospital. For a patient arriving directly at an EVT-capable center at onset time $t_a$, needle time is $t_a + T_{\text{DTN}}$ and puncture time $t_a + T_{\text{DTP}}$. At a thrombolysis-capable center, needle time is $t_a + T_{\text{DTN}}$ and, if the patient is transferred to an EVT-capable center at road time $\tau'$, puncture time is $t_a + T_{\text{DIDO}} + \tau' + T_{\text{DTP}}^{\text{tr}}$; which center, and whether to transfer at all, is the second decision. At a non-stroke-center hospital both treatments follow a transfer. Inter-facility transfers are not distance-limited; a pair is unreachable only if the road network has no route. For comparison, Holodinsky \emph{et al.}~\cite{holodinsky2018} assume door-to-needle 30, door-in-door-out 50 and door-to-puncture 60 (direct) or 30 (transfer) minutes in their base case, and 60 / 120 / 60--90 minutes in their slower scenarios; our base case lies between the two and is anchored to current registry medians and national targets.

\begin{table}[t]
\centering
\caption{In-hospital intervals in the base case (minutes) and the ranges examined in sensitivity analysis.}
\label{tab:intervals}
\resizebox{\columnwidth}{!}{\begin{tabular}{@{}lccl@{}}
\toprule
Interval & Base & Range & Source \\
\midrule
Door-to-needle $T_{\text{DTN}}$ & 45 & 30--60 & Target: Stroke III~\cite{targetstroke3} \\
Door-to-puncture, direct $T_{\text{DTP}}$ & 90 & 60--120 & Target: Stroke III~\cite{targetstroke3} \\
Door-to-puncture, transfer $T_{\text{DTP}}^{\text{tr}}$ & 60 & 45--90 & Target: Stroke III~\cite{targetstroke3} \\
Door-in-door-out $T_{\text{DIDO}}$ & 121 & 60--175 & GWTG-Stroke median, IQR 89--175~\cite{royan2026} \\
\bottomrule
\end{tabular}}
\end{table}


\subsection{Reward}
The reward is paid once, on the in-hospital decision that ends the episode, and is zero on the first leg. Its value is the probability of a good outcome, defined as an excellent outcome at 90 days (modified Rankin Scale 0--1), following Holodinsky \emph{et al.}~\cite{holodinsky2018}, evaluated at the needle and puncture times that the completed pathway implies: the first hospital $h_1$ reached at clock $t_1$, and the final hospital $h_2$ reached at $t_2$ ($h_2 = h_1$ if the patient stays). Thrombolysis is given at the first hospital on the pathway able to give it, so $t_n = t_1 + T_{\text{DTN}}$ if $h_1$ is thrombolysis- or EVT-capable and $t_n = t_2 + T_{\text{DTN}}$ if only $h_2$ is; thrombectomy requires an EVT-capable $h_2$, with $t_p = t_1 + T_{\text{DTP}}$ after direct arrival and $t_p = t_2 + T_{\text{DTP}}^{\text{tr}}$ after a transfer. For ischemic subtypes the outcome declines with onset-to-treatment time; for hemorrhage and mimics it is constant:
{\small\begin{align}
P_{\text{LVO,IVT}}(t_n) &= \max\{0.1328,\; 0.2359 - 4\!\times\!10^{-4} t_n + 2\!\times\!10^{-7} t_n^2\},\\
P_{\text{LVO,EVT}}(t_p) &= \max\{0.129,\; 0.3394 - 2\!\times\!10^{-4} t_p + 4\!\times\!10^{-8} t_p^2\},\\
P_{\text{LVO}} &= P_{\text{LVO,IVT}} + (1 - P_{\text{LVO,IVT}})\,P_{\text{LVO,EVT}},\\
P_{\text{nLVO}}(t_n) &= \max\{0.4622,\; 0.6343 - 5\!\times\!10^{-4} t_n - 5\!\times\!10^{-8} t_n^2\},\\
P_{\text{ICH}} &= 0.24, \qquad P_{\text{mimic}} = 0.90,
\end{align}}
where a treatment that the pathway does not provide contributes its untreated floor ($P_{\text{LVO,EVT}} = 0$ when no thrombectomy is performed), and $t_n, t_p \geq 270$ minutes yield the floor values. Because the reward is paid after imaging, it is always evaluated with the subtype known; subtype uncertainty enters only through the planner's expectation at the first decision (Section~\ref{sec:method}). Paying the reward once, rather than scoring each leg separately, is what makes the planner's objective and the reported outcome the same quantity.
